% ---------------------------------------------------------------- 摘要
\section*{摘要}
\addcontentsline{toc}{section}{摘要}

\begin{keybox}[报告要点]
\small
本报告以\textbf{申万行业分类（2021 版）}为统一框架，分为两大部分：
第一部分从供需关系出发分析农业各子行业的周期性，以\hl{猪周期为基准}
测算跨品种、跨行业的相关性与传导时滞；第二部分系统梳理
申万农林牧渔行业\textbf{全部 \num{104} 家 A 股上市公司}近三年的经营与融资画像。

\textbf{第一部分的主要发现}：
（1）本轮农业下行（\num{2024}---\num{2025}）源于\hl{供给恢复而非需求萎缩}——
玉米、小麦、稻米同时丰产，生猪出栏在产能未出清时继续增长，
白羽肉鸡祖代更新量创历史新高；
（2）品种周期长度由\hl{产能调整的生物学时间}决定，从蛋鸡的 \num{6}---\num{12} 个月
到肉牛的 \num{24}---\num{44} 个月以上，谱系跨越 \num{7} 倍；
（3）在统一口径（月度收益率、同一时间窗口）下，
\hl{跨品种相关性的绝对值普遍低于 \num{0.25}}——玉米 \num{+0.09}、豆粕 \num{+0.13}，
仅鸡蛋与生猪周期同源，“猪价带动一切农产品”并非统计事实；
（4）截至 \num{2026} 年 9 月，生猪自繁自养已连续六个月处于亏损区间
（猪粮比价月度均值 \num{3.96}---\num{4.47}）、能繁母猪存栏同比转负，
而肉牛价格、蛋鸡存栏与水产投苗已率先反转，\hl{农业周期正处于
“猪链磨底、禽链与牛链上行”的分化阶段}。

\textbf{第二部分的主要发现}：
（1）\num{104} 家公司中 \num{34} 家 2025 年亏损、\num{29} 家 ROE 为负，
ROE 中位数仅 \num{2.64}\%，而牧原、温氏两家合计利润 \num{207.5} 亿元，
\hl{等于全部养殖公司的利润总额}；
（2）子行业分化鲜明：动保（ROE 中位数 \num{7.10}\%）、宠物食品（平均 ROE \num{11.5}\%）
为抗周期方向，饲料（ROE 中位数 \num{-2.46}\%、亏损面 \num{53.3}\%）为最受损环节；
（3）\num{18} 家公司同时处于亏损与资产负债率 $>$ \num{65}\% 状态
（其中 \num{7} 家 $>$ \num{80}\%），\num{31} 家连续三年未分红；
据财务约束与公告行为推断，第一档“补血型”融资需求总规模在 \num{300}---\num{500} 亿元量级，
\hl{当前农业上市公司的资本运作主题是修复资产负债表，而非扩大产能}。
\end{keybox}

\vspace{6pt}
\noindent\textbf{关键词}\quad 农业周期；猪周期；跨品种相关性；申万行业分类；
上市公司画像；融资需求
\clearpage
